\begin{center}
	\large{\textbf{FRACTIONAL SPACE QUANTUM HEAT ENGINE PERFORMANCE REVIEW BASED ON 1 DIMENSION INFINITE POTENTIAL}}
\end{center}
\begin{minipage}{5cm}
\begin{align*}
&\text{Name} \qquad &&: \text{MUHAMMAD SULTHAN ZACKY RAIHAN PUTERA HASAN} \\
&\text{NRP} \qquad &&: \text{5001201046} \\
&\text{Department} \qquad &&: \text{Physics} \\
&\text{Advisor}\qquad &&: \text{Dr. Heru Sukamto, M.Si}\\
\end{align*}
\end{minipage} \\
\textbf{Abstract}
\par The Schrödinger equation can be approached in multiple ways, one of which incorporates fractional calculus, introducing a new parameter called the fractional parameter \(\alpha\). The eigenenergy equation obtained by solving the Schrödinger equation is now also affected by this new parameter. This eigenenergy equation is then applied to one branch of quantum mechanics, the quantum heat engine, to investigate how the fractional parameter affects the efficiency and power output of the engine. This study reviews three cycles: Carnot, Otto, and Brayton, using a particle trapped in one dimension infinite potential well as the working substance. It is found that the efficiency can be reduced to only a function of compression ratio \(r\) and fractional parameter \(\alpha\). The efficiency of all cycles is also found to have the same form, while the power output varies between them. It is evident that the fractional parameter \(\alpha\) can significantly impact both the efficiency and power output of a quantum heat engine. \\
\begin{minipage}{1cm}
	\begin{align*}
	\text{\textbf{Keywords}}:\hspace{2mm}&\text{\textit{\textbf{Efficiency, Fractional Parameter, Quantum Heat Engine}}} 
	\end{align*}
\end{minipage}

\newpage
\thispagestyle{empty}
\vspace*{\fill}
    \begin{center}
	Halaman ini sengaja dikosongkan
    \end{center}
\vspace*{\fill}
\pagebreak